% Ad Familiares 15.13 (ad-familiares-15-13)
% Source: https://raw.githubusercontent.com/PerseusDL/canonical-latinLit/master/data/phi0474/phi056/phi0474.phi056.perseus-lat1.xml
% Fetched by scripts/fetch_latin.py

% Dateline (Perseus): Scr. eodem loco et tempore quo ep. iv et x

\ciceroLetterOpener{M. CICERO IMP. S. D. L. PAVLLO COS.}

\ciceroSection{1} maxime mihi fuit optatum Romae esse tecum multas ob causas sed praecipue ut et in petendo et in gerendo consulatu meum tibi debitum studium perspicere posses. ac petitionis quidem tuae ratio mihi semper fuit explorata, sed tamen navare operam volebam ; in consulatu vero cupio equidem te minus habere negoti, sed moleste fero me consulem tuum studium adulescentis perspexisse, te meum, cum id aetatis sim, perspicere non posse.

\ciceroSection{2} sed ita fato nescio quo contigisse arbitror ut tibi ad me ornandum semper detur facultas, mihi ad remunerandum nihil suppetat praeter voluntatem. ornasti consulatum, ornasti reditum meum . incidit meum tempus rerum gerendarum in ipsum consulatum tuum. itaque cum et tua summa amplitudo et dignitas et meus magnus honos magnaque existimatio postulare videatur ut a te plurimis verbis contendam ac petam ut quam honorificentissimum senatus consultum de meis rebus gestis faciendum cures, non audeo vehementer a te contendere, ne aut ipse tuae perpetuae consuetudinis erga me oblitus esse videar aut te oblitum putem.

\ciceroSection{3} qua re, ut te velle arbitror, ita faciam atque ab eo, quem omnes gentes sciunt de me optime meritum, breviter petam. si alii consules essent, ad te potissimum, Paulle, mitterem ut eos mihi quam amicissimos redderes ; nunc cum tua summa potestas summaque auctoritas notaque omnibus nostra necessitudo sit, vehementer te rogo ut et quam honorificentissime cures decernendum de meis rebus gestis et quam celerrime. dignas res esse honore et gratulatione cognosces ex iis litteris, quas ad te et conlegam et senatum publice misi. omniumque mearum reliquarum rerum maximeque existimationis meae procurationem susceptam velim .habeas, in primisque tibi curae sit, quod abs te superioribus quoque litteris petivi, ne mihi tempus prorogetur. cupio te consulem videre omniaque quae spero cum absens tum etiam praesens te consule adsequi
